\documentclass[aspectratio=169]{beamer}
\usetheme{Madrid}
\usepackage{amsmath}
\usepackage{amsfonts}
\usepackage{amssymb}
\usepackage{bm}

\title{Theoretical Background}
\subtitle{Kitaev Chain, Topological Phases, and Majorana Edge Modes}
\author{}
\date{April 13, 2026}

\begin{document}

\begin{frame}
\titlepage
\end{frame}

\begin{frame}{Project Context}
\begin{itemize}
\item Goal: study the topological phase transition of the one-dimensional Kitaev chain.
\item Main bridge: connect fermionic many-body physics to qubit-based simulation methods.
\item Core tasks later in the project:
\begin{itemize}
\item derive and analyze the model,
\item map fermions to qubits,
\item prepare ground states and measure topological observables,
\item test the stability of signatures under noise.
\end{itemize}
\end{itemize}
\end{frame}

\begin{frame}{Why This Model Matters}
\begin{itemize}
\item The Kitaev chain is the minimal model for one-dimensional topological superconductivity.
\item It exhibits two distinct phases:
\begin{itemize}
\item a trivial phase with no protected edge physics,
\item a topological phase with Majorana zero modes localized at the boundaries.
\end{itemize}
\item It is simple enough for exact analysis and small-scale simulation.
\item It captures the main ideas behind edge modes, bulk-boundary correspondence, and topological protection.
\end{itemize}
\end{frame}

\begin{frame}{Kitaev Chain Hamiltonian}
\[
H =
-\mu \sum_{j=1}^{N} \left(c_j^\dagger c_j - \frac{1}{2}\right)
- t \sum_{j=1}^{N-1} \left(c_j^\dagger c_{j+1} + c_{j+1}^\dagger c_j\right)
+ \Delta \sum_{j=1}^{N-1} \left(c_j c_{j+1} + c_{j+1}^\dagger c_j^\dagger\right)
\]

\begin{itemize}
\item \(c_j^\dagger, c_j\): spinless fermion creation and annihilation operators.
\item \(t\): nearest-neighbor hopping amplitude.
\item \(\mu\): chemical potential.
\item \(\Delta\): superconducting pairing amplitude.
\end{itemize}
\end{frame}

\begin{frame}{Physical Interpretation}
\begin{itemize}
\item The hopping term delocalizes fermions along the chain.
\item The chemical potential controls filling and tunes the phase transition.
\item The pairing term breaks particle-number conservation but preserves fermion parity.
\item The Hamiltonian is quadratic, so it can be solved exactly by Bogoliubov methods.
\end{itemize}
\[
\text{number not conserved} \qquad \Longrightarrow \qquad \text{Nambu/BdG description}
\]
\end{frame}

\begin{frame}{Momentum-Space BdG Form}
\[
\Psi_k =
\begin{pmatrix}
c_k \\
c_{-k}^\dagger
\end{pmatrix},
\qquad
H = \frac{1}{2}\sum_k \Psi_k^\dagger \mathcal{H}(k)\Psi_k
\]
\[
\mathcal{H}(k) =
\left(-2t\cos k - \mu\right)\tau_z + 2\Delta \sin k \, \tau_y
\]
\[
E(k) = \pm \sqrt{\left(2t\cos k + \mu\right)^2 + 4\Delta^2 \sin^2 k}
\]

\begin{itemize}
\item The spectrum is gapped except when the bulk gap closes.
\item Gap closings mark topological phase transitions.
\end{itemize}
\end{frame}

\begin{frame}{Phase Diagram}
\begin{itemize}
\item The bulk gap closes at
\[
\mu = \pm 2t
\]
\item For \(\Delta \neq 0\):
\begin{itemize}
\item \(|\mu| < 2|t|\): topological superconducting phase,
\item \(|\mu| > 2|t|\): trivial phase.
\end{itemize}
\item The topological phase supports boundary Majorana zero modes for open chains.
\end{itemize}
\[
\text{bulk gap closing} \Longleftrightarrow \text{change of topological sector}
\]
\end{frame}

\begin{frame}{Majorana Representation}
\[
c_j = \frac{1}{2}\left(a_j + i b_j\right),
\qquad
c_j^\dagger = \frac{1}{2}\left(a_j - i b_j\right)
\]
\[
a_j^\dagger = a_j,\qquad b_j^\dagger = b_j,\qquad
\{a_j,a_\ell\}=2\delta_{j\ell},\qquad
\{b_j,b_\ell\}=2\delta_{j\ell}
\]

\begin{itemize}
\item Each fermionic mode splits into two Majorana operators.
\item In the trivial phase, Majoranas pair on the same site.
\item In the topological phase, Majoranas pair between neighboring sites, leaving unpaired edge modes.
\end{itemize}
\end{frame}

\begin{frame}{Special Point and Edge Modes}
\begin{itemize}
\item At the solvable point \(t=\Delta\) and \(\mu=0\), the structure is especially transparent.
\item The Hamiltonian pairs \(b_j\) with \(a_{j+1}\) in the bulk.
\item The operators \(a_1\) and \(b_N\) remain unpaired for an open chain.
\item These boundary operators form a nonlocal fermionic zero mode.
\end{itemize}
\[
f_{\text{edge}} = \frac{1}{2}\left(a_1 + i b_N\right)
\]
\begin{itemize}
\item This produces a near-degenerate ground-state manifold distinguished by fermion parity.
\end{itemize}
\end{frame}

\begin{frame}{Finite Chains and Observable Signatures}
\begin{itemize}
\item In a finite chain, the two edge Majoranas overlap weakly.
\item This overlap causes an exponentially small energy splitting.
\item Practical signatures of the topological regime include:
\begin{itemize}
\item a bulk gap that closes and reopens,
\item near-zero-energy edge states,
\item boundary-localized correlations,
\item parity-sensitive ground-state structure.
\end{itemize}
\item These are the quantities that simulation and measurement protocols must recover.
\end{itemize}
\end{frame}

\begin{frame}{Fermion-to-Qubit Mapping}
\[
c_j =
\left(\prod_{\ell<j} Z_\ell\right)\frac{X_j - iY_j}{2},
\qquad
c_j^\dagger =
\left(\prod_{\ell<j} Z_\ell\right)\frac{X_j + iY_j}{2}
\]

\begin{itemize}
\item The Jordan-Wigner map converts fermionic operators into Pauli strings.
\item This provides the bridge from many-body theory to circuit simulation.
\item After the mapping, the model can be implemented with qubit Hamiltonians and measured with standard quantum-information tools.
\end{itemize}
\end{frame}

\begin{frame}{What This Background Enables}
\begin{itemize}
\item Exact diagonalization of small systems.
\item Free-fermion calculations of spectra and correlation functions.
\item Circuit-level emulation of state preparation and measurement.
\item Later analysis of robustness under noise and imperfect operations.
\end{itemize}
\[
\text{theory} \;\rightarrow\; \text{encoding} \;\rightarrow\; \text{simulation} \;\rightarrow\; \text{measurement}
\]
\end{frame}

\begin{frame}{Takeaway}
\begin{itemize}
\item The Kitaev chain is the central theoretical model for the project.
\item Its phase transition is controlled by a bulk gap closing at \(\mu=\pm 2t\).
\item The topological phase supports Majorana edge modes in open chains.
\item The Jordan-Wigner transformation makes the model accessible to qubit-based simulation.
\item This is the theoretical foundation for all later numerical and circuit work.
\end{itemize}
\end{frame}

\end{document}
